% ID: SEL-D0069-OB-E02
% ESTADO: validada. S069.
\textbf{a)} Al completar el cuadrado, $C(t)=45-0{,}2(t-10)^2$. El máximo se alcanza en $t=10$, dentro del dominio, que corresponde al año $\boxed{2010}$.

\textbf{b)} $C(t)=0$ equivale a $(t-10)^2=225$, de donde $t=-5$ o $t=25$. Sólo $25$ pertenece al intervalo del modelo. El nivel cero corresponde al año $\boxed{2025}$.

\textbf{c)} $C'(t)=4-0{,}4t$, por lo que la pendiente en $8$ es $\boxed{0{,}8}$. Es positiva: en torno a ese instante el nivel de contaminación está aumentando, a razón instantánea de $0{,}8$ unidades de nivel por año.
